% YouthTwin manuscript, IEEE journal format (IEEEtran).
% Every number is a macro from ../tables/ieee_values.tex or a table from ../tables/, generated by
% scripts/build_paper.py from runs. Do not type numbers here (results guard, working rule 4).
% Build: make paper-pdf
\documentclass[journal]{IEEEtran}
\usepackage[T1]{fontenc}
\usepackage{amsmath,amssymb}
\usepackage{graphicx}
\usepackage{booktabs}
\usepackage{xcolor}
\usepackage{cite}
\usepackage{url}
\usepackage[hidelinks]{hyperref}
\graphicspath{{../figures/}}
\input{../tables/ieee_values.tex}
\providecommand{\PENDINGFIG}{\textcolor{red}{[figure pending]}}

\begin{document}

\title{An Uncertainty-Aware Growth Twin for Maturity-Conditioned Injury-Risk Prediction in Youth Athletes: Validation on Real Growth and Injury Data and a Calibrated Simulation Study}

% TODO(human): authors, affiliations, funding and ethics statements.
\author{Author~Names~Withheld%
\thanks{Code: this repository; every number is traceable to a run record.}}

\markboth{Manuscript draft}{Growth twin for youth injury risk}

\maketitle

\begin{abstract}
Growth-related injuries in youth sport cluster around the adolescent growth spurt, yet maturity
is usually entered into injury models as a point estimate from an equation, without
uncertainty, and published injury-prediction models rarely report calibration. We present a
digital twin that (i) estimates each athlete's age at peak height velocity (PHV) with a
hierarchical Bayesian growth-curve model fitted to the sparse measurements of a single season,
(ii) propagates that uncertainty into a discrete-time injury-hazard model, and (iii) takes
movement features from smartphone video. Each layer is tested on the best available evidence.
On real longitudinal growth records of \BerkChildren{} children (Berkeley Child Guidance
Study), given one season of three measurements per child, the twin estimated age at PHV with a
mean absolute error of \BerkTwinMAE{}~years against \BerkMooreMAE{}~years for the Moore
equation, with \BerkCoverage{} coverage of nominal \SpecInterval{} intervals. On real injury
data from \RunRunners{} competitive runners, athlete-grouped evaluation gave an AUC of
\RunAucLogit{} (\RunAucLogitCI{}), and allowing an athlete's days on both sides of a split
inflated AUC by up to \RunInflBoost{}. A simulation study calibrated to published incidence,
growth-velocity effects and smartphone measurement error shows that knee flexion meets a
\SpecRmseLimit{}$^\circ$ agreement limit at every maturity stage, whereas frontal-plane knee
angles do not reliably do so before PHV, and quantifies the cohort size a prospective study
needs. All simulated results are labelled as such and are not findings about athletes.
\end{abstract}

\begin{IEEEkeywords}
Digital twin, growth curve, peak height velocity, Bayesian hierarchical model, injury
prediction, calibration, pose estimation, youth sport.
\end{IEEEkeywords}

\section{Introduction}
\IEEEPARstart{G}{rowth-related} injuries, such as apophysitis at the knee and heel and
physeal stress injuries, are frequent in youth sport and peak around the adolescent growth
spurt. In English academy football, apophyseal injuries were recorded at a rate of about
\SpecIncidence{} of player-seasons \cite{oxendale2025}, and faster growth is consistently
associated with higher injury risk \cite{hall2025}. Prediction of these injuries is therefore an
attractive target, but the field's record is poor: of 204 published musculoskeletal
injury-prediction models, almost all were at high or unclear risk of bias, few reported
calibration, and none was externally validated \cite{bullock2022}.

Two weaknesses recur. First, maturity enters models as a point estimate from an anthropometric
equation \cite{mirwald2002,moore2015}. These equations regress towards the mean, are least
accurate far from PHV, and carry no statement of uncertainty, so a model cannot distinguish an
athlete whose maturity is well determined from one whose maturity is a guess. Second,
evaluation often lets correlated observations of the same athlete fall on both sides of a
train--test split, which flatters performance.

We address both with a \emph{growth twin}: a hierarchical Bayesian growth-curve model in the
SITAR family \cite{cole2010sitar} that estimates each athlete's growth state, including its
uncertainty, from the few measurements of one season, and passes draws of that state to an
injury-hazard model. Movement is measured from smartphone video with a 2D pose estimator
\cite{bazarevsky2020blazepose}, and a conditioned neural encoder lets the growth state modulate
how movement is read \cite{perez2018film}.

No public dataset combines youth growth, smartphone video and diagnosed growth-related injury.
We therefore test each layer on the strongest evidence available: the growth twin on real
longitudinal growth records; the evaluation discipline on real injury outcomes; and the layers
without public data in a simulation study calibrated to published values. Our contributions
are: (1) a single-season growth twin validated against real growth records, with honest
uncertainty; (2) an evaluation protocol shown on real injuries to expose leakage and
miscalibration; (3) a calibrated simulation that states what smartphone kinematics and a
growth-aware risk model can deliver, and how large a prospective cohort must be.

\section{Methods}

\subsection{Growth twin}
For athlete $i$ with stature $y_{ij}$ at age $t_{ij}$,
\begin{equation}
y_{ij} = \mu + \alpha_i + \delta_i \,(t_{ij} - \bar t) + h\!\big((t_{ij}-\beta_i)\,e^{\gamma_i}\big) + \varepsilon_{ij},
\end{equation}
where $h$ is a natural cubic spline shared by all athletes of one sex, $\alpha_i$ shifts size,
$\beta_i$ shifts timing, $\gamma_i$ scales the tempo of the spurt, $\delta_i$ absorbs a mismatch
in velocity between the athlete and the shared curve, and $\varepsilon_{ij}$ is measurement
error. The prior for $\beta_i$ is centred on the athlete's Moore estimate \cite{moore2015},
with a cohort-level offset and spread learned from the data. With one season of data, timing
is weakly identified and its posterior can be two-valued; sampling $\beta_i$ directly then
fails. We therefore integrate $\alpha_i$ and $\delta_i$ analytically and $\beta_i$ and
$\gamma_i$ numerically on grids, sample only the shared parameters with NUTS, and draw the
athlete effects exactly given each shared draw. Age at PHV is the age of the highest interior
peak of the athlete's velocity curve. Full specification and priors are in the code and the
generated methods text.

\subsection{Real growth data}
The Berkeley Child Guidance Study \cite{tuddenham1954}, distributed with the \texttt{sitar}
package \cite{sitarpkg}, follows \BerkChildren{} children born in Berkeley, California, in
1928--29, with stature, stem (sitting) height and mass every six months from 8 to 21 years.
It was not used to develop or calibrate the Moore equations. A reference age at PHV was
obtained for each child from the full record by fitting Preece--Baines model 1
\cite{preece1978}, checked against the velocity peak of a smoothing spline; children whose
record did not span the spurt, whose fit was poor, or whose two estimates disagreed were not
scored (\BerkIncluded{} children scored; \BerkStrict{} under a stricter agreement rule). The
estimators were then given one season per child: three consecutive six-monthly measurements
starting at a random age between \SpecWindowLow{} and \SpecWindowHigh{} years, drawn \BerkReplicates{} times. The twin was
fitted per sex to every child's season; Moore \cite{moore2015} and, for boys, Mirwald
\cite{mirwald2002} were computed from the same measurements.

\subsection{Real injury data}
To test the evaluation protocol on real outcomes we used the open dataset of
L\"ovdal et~al. \cite{lovdal2021,lovdal2021data}: \RunRunners{} competitive runners,
\RunRows{} event days each with three weeks of aggregated training load, and \RunInjuries{}
injury days. Models were a load-only logistic model, an elastic-net logistic model on all
inputs and a gradient-boosted model. Week-to-week load ratios, which divide by zero in rest
weeks, were replaced by bounded log ratios computed from each row alone. Three schemes were
compared: five-fold cross-validation grouped by runner (primary); a temporal hold-out of the
last quarter of calendar time; and row-level five-fold cross-validation, which lets a runner's
nearly identical days fall in both training and test sets and is reported only to measure the
resulting inflation. All preprocessing was fitted on training folds only.

\subsection{Simulation study}
\emph{Cohort.} Synthetic athletes were generated with Preece--Baines curves whose shape and
spread were taken from the Berkeley fits, with mean age at PHV set to Indian reference values
(\SpecFemaleAPHV{} years in girls, \SpecMaleAPHV{} in boys \cite{khadilkar2019}), visits every
six to eight weeks, and persistent week-to-week training load. Weekly cause-specific injury
hazards were known by construction. The growth-related hazard rose by \SpecRiskIncrease{} per
\SpecRiskGrowth{}~cm/year of lower-limb growth \cite{hall2025}, and its intercept was solved so
that \SpecIncidence{} of athletes sustained a growth-related injury per season
\cite{oxendale2025}. Effects of a landing-load proxy, of its interaction with growth velocity,
and of training hours have no published estimates and are assumptions.

\emph{Smartphone versus motion capture.} Each simulated drop jump was observed twice: by a
marker-like reference and by a smartphone-like estimator with per-participant and per-trial
landmark offsets and frame noise. Both observations passed through the full pose pipeline
(gap filling, zero-lag low-pass filtering, event detection, angle definitions). The landmark
offset was calibrated so that peak knee flexion had a mean absolute error of
\SpecTargetMAE{}$^\circ$ in adult-sized participants, the value reported for a smartphone
system against Vicon \cite{barzyk2024} (\ValOffset{}~px). Errors were expressed in pixels at
adult size and scaled by body size, so any loss of agreement in smaller, less mature children is
an output of the simulation. The planned study of \SpecValN{} participants (equal pre-, circa-
and post-PHV groups) was simulated once and repeated \ValStudies{} times. Agreement used
Bland--Altman limits \cite{bland1986}, the two-way random-effects absolute-agreement intraclass correlation (ICC) and RMSE against the pre-specified
\SpecRmseLimit{}$^\circ$ limit.

\emph{Risk models.} In a simulated cohort of \SimAthletes{} athletes followed for one season,
an athlete-week table was built with the growth-twin state, observed growth velocities,
screening features and exposure; the outcome was a growth-related injury in the next four
weeks. The confirmatory model used as many pre-specified predictors as the Riley criteria
\cite{riley2020} allow for the achieved events (\SimConfPredictors). Exploratory models were a
penalised logistic hazard, a boosted hazard and neural hazards with a temporal convolutional
movement encoder conditioned on the growth state, compared in a nine-configuration ablation
with deep ensembles \cite{lakshminarayanan2017}. Evaluation used athlete-grouped five-fold
cross-validation with athlete-level bootstrap intervals, calibration intercept and slope
\cite{vancalster2019}, and decision curves \cite{vickers2006dca}. A cohort-size study repeated
the tabular evaluation on independent simulated cohorts of increasing size.

\emph{Reporting.} The study follows TRIPOD+AI \cite{collins2024tripodai}. Hypotheses, the
primary outcome and the ablation set were pre-registered before any real outcome data were
analysed.

\section{Results}

\subsection{Growth twin on real growth records}
Mean reference age at PHV was \BerkGirlsAPHV{} years in girls and \BerkBoysAPHV{} in boys.
Table~\ref{tab:growth} compares the estimators over \BerkWindows{} child-seasons. Overall, the
twin's mean absolute error was \BerkTwinMAE{}~years against \BerkMooreMAE{}~years for Moore
and \BerkMirwaldMAE{}~years for Mirwald (boys), with mean biases of \BerkTwinBias{} and
\BerkMooreBias{}~years. Nominal \SpecInterval{} intervals covered the reference in \BerkCoverage{}
of child-seasons: \BerkCoverageNear{} for seasons within three years of PHV and
\BerkCoverageFar{} further away. Within three years of PHV the errors were \BerkTwinMAENear{}
and \BerkMooreMAENear{}~years. Restricting scoring to the \BerkStrict{} children whose reference
methods agreed closely gave \BerkStrictTwinMAE{} and \BerkStrictMooreMAE{}~years and
\BerkStrictCoverage{} coverage. Across fits the largest R-hat was \BerkRhat{} and the smallest
bulk effective sample size \BerkESS{}.

\begin{table*}[!t]
\caption{Age at PHV from one season of measurements, real Berkeley data. Mean absolute error and
bias in years; coverage of nominal \SpecInterval{} intervals; mean interval width in years.
Bands: years from PHV at mid-season.}
\label{tab:growth}
\centering
\input{../tables/ieee_tab_growth.tex}
\end{table*}

\begin{figure}[!t]
\centering
\IfFileExists{../figures/ieee_fig_growth.png}{\includegraphics[width=\columnwidth]{ieee_fig_growth.png}}{\fbox{\PENDINGFIG}}
\caption{Estimated versus reference age at PHV on real data (one window draw): growth twin with
intervals (left) and Moore equation (right).}
\label{fig:growth}
\end{figure}

\subsection{Evaluation on real injury outcomes}
With runner-grouped cross-validation (Table~\ref{tab:runners}), the load-only model did not
discriminate (AUC \RunAucLoad{}, \RunAucLoadCI{}), whereas the penalised logistic model reached
\RunAucLogit{} (\RunAucLogitCI{}); its paired AUC gain over the load-only model was
\RunDiffLogit{} (\RunDiffLogitCI{}). The temporal hold-out gave a similar AUC (\RunTempAucLogit{})
and a calibration slope of \RunTempSlopeLogit{}, while under grouped cross-validation the slopes
were below one (\RunSlopeLogit{} and \RunSlopeBoost{}), indicating predictions that are too
extreme. Letting a runner's days fall on both sides of the split raised AUC by \RunInflLoad{},
\RunInflLogit{} and \RunInflBoost{} for the three models, each with an interval excluding zero.
The dataset's authors reported an AUC of \CitedPublishedAuc{} for this task \cite{lovdal2021};
we do not attribute the difference to any particular cause.

\begin{table}[!t]
\caption{Real injury outcomes (competitive runners). Athlete-level bootstrap intervals.}
\label{tab:runners}
\centering
\resizebox{\columnwidth}{!}{\input{../tables/ieee_tab_runners.tex}}
\end{table}

\subsection{Simulation study: smartphone versus motion capture}
\emph{All results in this and the following subsections are from synthetic data.}
Peak knee flexion met the \SpecRmseLimit{}$^\circ$ limit in at least \ValKneePassMin{} of
simulated studies in every maturity group (mean RMSE \ValKneeRMSEPre{}$^\circ$ before PHV,
\ValKneeRMSEPost{}$^\circ$ after; Table~\ref{tab:validity}). The knee frontal-plane projection
angle at peak flexion was less accurate and degraded in smaller children: mean RMSE
\ValFppaRMSEPre{}$^\circ$ before, \ValFppaRMSECirca{}$^\circ$ around and \ValFppaRMSEPost{}$^\circ$
after PHV, so that a study of \SpecValN{} participants would conclude RMSE below the limit in
\ValFppaPassPre{}, \ValFppaPassCirca{} and \ValFppaPassPost{} of repetitions respectively.

\begin{table*}[!t]
\caption{Simulated smartphone versus motion capture (SYNTHETIC). Bias, limits of agreement and
ICC from one study of \SpecValN{} participants; RMSE mean and \SpecCI{} range, and share of repeated
studies with RMSE below the \SpecRmseLimit{}$^\circ$ limit, from \ValStudies{} repetitions.
Angles in degrees.}
\label{tab:validity}
\centering
\input{../tables/ieee_tab_validity.tex}
\end{table*}

\subsection{Simulation study: risk models}
In the simulated cohort of \SimAthletes{} athletes, \SimEventWeeks{} athlete-weeks were followed
by a growth-related injury within four weeks. Table~\ref{tab:simrisk} reports out-of-fold
performance. The confirmatory model reached an AUC of \SimConfAuc{} (\SimConfAucCI{}) with
calibration slope \SimConfSlope{}; the exploratory penalised hazard reached \SimLogitAuc{}
(\SimLogitAucCI{}). The growth twin in this cohort estimated age at PHV with an error of
\SimTwinMAE{} against \SimMooreMAE{}~years for Moore, with \SimCoverage{} interval coverage.

\begin{table}[!t]
\caption{Simulated risk models (SYNTHETIC), growth-related injury within four weeks.}
\label{tab:simrisk}
\centering
\resizebox{\columnwidth}{!}{\input{../tables/ieee_tab_simrisk.tex}}
\end{table}

\subsection{Simulation study: conditioning ablation}
Table~\ref{tab:ablation} lists the nine pre-registered configurations. The paired AUC
difference of FiLM conditioning was \AblFilmVsConcat{} (\AblFilmVsConcatCI{}) against input
concatenation, \AblFilmVsLate{} (\AblFilmVsLateCI{}) against late fusion and \AblFilmVsNone{}
(\AblFilmVsNoneCI{}) against a movement encoder without growth state. Nominal \SpecInterval{}
risk intervals from growth-draw propagation and ensembling contained the true simulated risk in
\AblCoverage{} of athlete-weeks.

\begin{table}[!t]
\caption{Conditioning ablation (SYNTHETIC), growth-related injury.}
\label{tab:ablation}
\centering
\resizebox{\columnwidth}{!}{\input{../tables/ieee_tab_ablation.tex}}
\end{table}

\subsection{Simulation study: cohort size}
With \SizeSmall{} athletes a season produced on average \SizeEventsSmall{} athletes with a
growth-related injury, and \SizeEventsLarge{} with \SizeLarge{} athletes (Table~\ref{tab:sizes},
Fig.~\ref{fig:sizes}). The penalised model's AUC was \SizeLogitAucSmall{} and
\SizeLogitAucLarge{} at these sizes against a ceiling of \SizeTrueAucLarge{} set by the true
risk, and its calibration slope moved from \SizeLogitSlopeSmall{} to \SizeLogitSlopeLarge{}.

\begin{table*}[!t]
\caption{Out-of-fold performance against cohort size (SYNTHETIC). Mean and range across
independent simulated cohorts.}
\label{tab:sizes}
\centering
\input{../tables/ieee_tab_sizes.tex}
\end{table*}

\begin{figure}[!t]
\centering
\IfFileExists{../figures/ieee_fig_sizes.png}{\includegraphics[width=\columnwidth]{ieee_fig_sizes.png}}{\fbox{\PENDINGFIG}}
\caption{Performance against cohort size (SYNTHETIC); labels give the mean number of athletes
with a growth-related injury.}
\label{fig:sizes}
\end{figure}

\section{Discussion}
\emph{Growth twin.} On real children the twin's estimate is compared directly with the
equations practitioners use, using the same single season of data. Its value lies as much in
the interval as in the point estimate: the interval tells a practitioner when the data do not
determine maturity timing, which an equation cannot. Coverage is the property to check before
its uncertainty is used downstream, and it should be judged separately near and far from PHV.

\emph{Evaluation.} The runners' data show on real outcomes why athlete-grouped splitting is
non-negotiable: leaking an athlete across folds raised AUC by up to \RunInflBoost{}, enough to
reverse the ranking of models whose AUCs differ by less than that. Calibration slopes below one under grouped
evaluation show that discrimination alone would overstate readiness for use.

\emph{Smartphone kinematics.} The simulation, anchored to a published knee-flexion error,
predicts that sagittal knee and hip angles are measured well enough at every maturity stage,
whereas frontal-plane knee angles are marginal and degrade in smaller children. A validation
study should therefore be stratified by maturity and powered for the frontal plane.

\emph{Limitations.} The real growth data come from a historical US cohort, and the real injury
data from adult runners; neither tests growth-related injury in youth. The simulated layers
depend on assumptions stated in the methods, chiefly the effects of landing load and its
interaction with growth, and the landmark-error model for the frontal plane. The Mirwald girls'
equation was not used because its coefficients have not been verified against the original
publication. The propagated risk intervals reflect growth-timing and ensemble uncertainty only.

\section{Conclusion}
A single-season growth twin can be validated on real growth records today, and the evaluation
discipline it needs can be demonstrated on real injury data. The full claim, that a growth-aware
model predicts growth-related injury from smartphone video, requires a prospective cohort whose
size the simulation study sets out; this work supplies the model, the protocol and the
pre-registration for it.

\section*{Data and Code Availability}
Berkeley data: \texttt{sitar} R package (CRAN). Runners' data: DataverseNL,
doi:10.34894/UWU9PV (CC0). Synthetic data are regenerated from fixed seeds. Code, configuration
and run records are in the project repository; every number in this paper is generated from a
run and traceable to its configuration, seed and commit.

\bibliographystyle{IEEEtran}
\bibliography{references}

\end{document}
